%======================================================================
% Section 7: Nonlocal and Scale-Dependent Continuum Models
%======================================================================
\section{Nonlocal and Scale-Dependent Continuum Models}
\label{sec:nonlocal}

The constitutive laws of Sections~\ref{sec:thermoelastic} and~\ref{sec:multiphysics} all share an assumption that is worth stating explicitly because it is the assumption that a substantial body of later work set out to remove: that the stress at a point is a function of the strain at that point alone. The assumption is not a modelling convenience; it is the statement that the continuum has no internal length scale. Once a material is known to possess one---because its microstructure has a characteristic dimension, or because experiments reveal size-dependent stiffness---the classical constitutive law cannot be correct at small scales, and a class of models has been built in which the constitutive response becomes either spatially non-local or gradient-dependent.

This section develops that class. The organising question is not which of the competing formulations is correct, but \emph{which physical mechanism each one encodes}, because the mechanisms predict measurably different dispersion and therefore admit different experimental tests.

\subsection{Physical motivation and constitutive foundations}
\label{sec:7.1}

The classical assumption fails in a specific and observable way. A beam of a given material, tested in bending, exhibits a stiffness that depends on the specimen thickness: thinner specimens are stiffer than the classical prediction, and the departure becomes significant below a material-dependent critical length. Two distinct physical mechanisms can produce this behaviour, and they are the root of the entire family of models considered here.

The first mechanism is \emph{constitutive}: the internal stress in the material depends on the neighbouring material state, because the mechanism that carries load---bonding, dislocation motion, stress concentration in a heterogeneous microstructure---acts over a finite distance. Eringen's formulation of nonlocal elasticity is the canonical statement of this mechanism \cite{Eringen1972,Eringen1974}, in which the stress at a point is obtained by averaging the local stress over a neighbourhood weighted by a kernel:

\begin{equation}
\sigma_{ij}(\bm{x}) = \int_\Omega \alpha(|\bm{x}' - \bm{x}|, \tau)\, \sigma_{ij}^{\mathrm{local}}(\bm{x}')\, d\Omega_{\bm{x}'},
\label{eq:7.1}
\end{equation}
where $\alpha$ is the attenuation kernel and $\tau = e_0 a/\ell$ is the nonlocal parameter, with $e_0$ a material constant, $a$ an internal characteristic length such as the lattice spacing, and $\ell$ an external characteristic length such as the specimen size. The nonlocal parameter is therefore not an arbitrary length: it is a ratio of an internal to an external scale, which means that the nonlocal correction grows in relative importance as the specimen is made smaller relative to the microstructure. This is the scaling property that distinguishes a size-dependent theory from an ad hoc correction.

The second mechanism is \emph{structural}: the material possesses additional internal degrees of freedom---microrotations, relative displacement of the constituents of a porous skeleton, dislocation density---whose motion must be carried by additional stress measures. Micropolar elasticity, formulated by Eringen as a linear theory of a continuum with independent microrotation fields \cite{Eringen1966a} and extended to the thermoelastic case by Tauchert, Claus and Ariman \cite{Eringen1968b}, is the earliest and most complete statement of this mechanism. Later work extended it to materials with voids, in which the skeleton and the pore fluid carry independent kinematic fields \cite{Nunziato1979,Cowin1983}, with a corresponding thermoelastic formulation \cite{Iesan1986}.

\subsubsection{Differential form}

For special kernel functions---notably the one-parameter exponential and the Helmholtz kernel---the integral equation of Eq.~\eqref{eq:7.1} reduces to a differential equation of higher order. For the Helmholtz kernel one obtains

\begin{equation}
(1 - \xi^2\nabla^2)\sigma_{ij} = \sigma_{ij}^{\mathrm{local}},
\label{eq:7.2}
\end{equation}
where $\xi = e_0 a$ is the nonlocal length scale parameter. The differential form is the one used in wave-propagation applications, because it is a partial differential equation and can therefore be treated with the transform and numerical machinery of Section~\ref{sec:methods}. It is worth recording that the differential and integral forms are not interchangeable in general: the integral relation with a given kernel and the differential operator of Eq.~\eqref{eq:7.2} agree only for that kernel, and the difference matters for the boundary conditions, which in the differential form become higher-order \cite{Polizzotto2001}.

\subsection{Strain-gradient and couple-stress theories}
\label{sec:7.2}

Strain-gradient theory extends classical elasticity by including higher-order gradients of strain in the constitutive relations:

\begin{equation}
\sigma_{ij} = C_{ijkl}\varepsilon_{kl} - \ell^2 C_{ijkl}\nabla^2\varepsilon_{kl},
\label{eq:7.3}
\end{equation}

where $\ell$ is a material length scale parameter. The physical reading of Eq.~\eqref{eq:7.3} is that the stress depends not only on the strain but on the curvature of the strain field, so that regions in which the strain varies rapidly---near a notch, a contact edge, or an interface---carry additional stress. The coefficient $\ell$ is a genuine material constant that must be measured; the identification procedures for it, together with the finite element implementations and the variational formulations that make the theory tractable, have been surveyed by Askes and Aifantis \cite{Askes2006,Askes2011}.

The unifying point of view that connects gradient elasticity to micropolar and couple-stress elasticity is that all three introduce internal degrees of freedom and a length scale, and differ in whether those degrees of freedom are constrained to coincide with the macroscopic deformation or are allowed to evolve independently. Aifantis's formulation, in which the internal length is interpreted as the scale of the microstructural heterogeneity that generates it, is the common root \cite{Askes2006}. It follows that a continuum model with internal length, in whichever of these forms, can be used to represent materials for which no explicit microstructural model is available---which is why the same length-scale parameter appears in treatments of metals, polymers, granular media, and biological tissue with broadly comparable values.

\subsection{Size-dependent wave behavior: dispersion from nonlocality}
\label{sec:7.3}

Nonlocality introduces dispersion: the wave speed depends on wavelength. The differential form of nonlocal elasticity was developed by Eringen \cite{Eringen1983}, and the size-dependent elastic effects at the nanoscale were quantified by Maranganti and Sharma \cite{Maranganti2007}, who established the characteristic length below which classical elasticity fails for a range of materials, and by Maranganti, Sharma and Sharma \cite{Maranganti2006}, who showed that the nanoscale nonlocal contribution to electromechanical coupling is generically non-zero even for materials that possess no piezoelectric tensor---a result directly relevant to the coupled theories of Section~\ref{sec:multiphysics}, because it means that a purely elastic size effect can generate an apparent electromechanical response in a nominally non-piezoelectric material.

For a nonlocal elastic bar, the dispersion relation is

\begin{equation}
c(k) = \frac{c_0}{\sqrt{1 + \xi^2 k^2}},
\label{eq:7.4}
\end{equation}

where $c_0$ is the classical wave speed, $k$ is the wave number, and $\xi$ is the nonlocal length scale. Long waves ($k \to 0$) approach the classical speed $c_0$; short waves are strongly slowed.

\subsubsection{The dispersion signature distinguishes the mechanisms}

It is worth emphasising what Eq.~\eqref{eq:7.4} does and does not say, because the corresponding statement in gradient elasticity differs in a way that is experimentally discriminating. The nonlocal dispersion relation \emph{saturates}: as $k \to \infty$ the wave speed tends to zero like $1/k$, so that the high-wave-number branch of the dispersion curve is bounded. A gradient theory produces a dispersion relation in which the squared wave speed decreases \emph{linearly} in $k^2$ over the range in which the gradient expansion is valid, and the two mechanisms therefore predict measurably different wave speeds at the same wave number even when they are fitted to agree in the long-wave limit. A dispersion measurement over a sufficiently broad range of $k$ can thus distinguish a nonlocal mechanism from a gradient mechanism, and the two are not interchangeable even though both produce the same qualitative signature of a size-dependent stiffness. This is one of the few places in the reviewed literature where two competing formulations of the same effect make quantitatively different and testable predictions, and it deserves more attention than it has received.

\subsection{Energy flux in nonlocal media}
\label{sec:7.4}

For nonlocal media, the classical divergence theorem must be modified because the stress at a point depends on strain throughout a neighbourhood. The mechanical power $\sigma_{ji}v_i$ at a point includes contributions from the velocity field in a neighbourhood, not only from the value at the point itself. The energy flux takes the form

\begin{equation}
J_i = \sigma_{ji}v_i + q_i + J_i^{\mathrm{nonlocal}},
\label{eq:7.5}
\end{equation}

where $J_i^{\mathrm{nonlocal}}$ is the nonlocal correction arising from the integral constitutive relation. Formally, the correction is obtained by writing the stress power in the integral form and applying the divergence theorem to the integral representation:

\begin{equation}
J_i^{\mathrm{nonlocal}} = \int_\Omega \alpha(|\bm{x}' - \bm{x}|)\, \sigma_{ji}^{\mathrm{local}}(\bm{x}') v_i(\bm{x}')\, d\Omega_{\bm{x}'} - \sigma_{ji}(\bm{x})v_i(\bm{x}).
\label{eq:7.6}
\end{equation}

Two remarks. First, the correction is not an optional refinement: without it, the energy balance of a nonlocal medium does not close, and a calculation that computes energy ratios using the classical flux alone will not satisfy the conservation identity of Section~\ref{sec:10.5}. Second, the derivation of Eq.~\eqref{eq:7.6} requires a modified divergence theorem for integral-type constitutive relations, and the rigorous treatment of that theorem together with the associated variational structure is a delicate matter \cite{Polizzotto2001}. The explicit nonlocal energy flux for the thermoelastic and multiphysics settings is not available in the reviewed corpus in a form that can be independently verified, and it is recorded as an open item rather than asserted.

\subsection{Interface conditions for nonlocal media}
\label{sec:7.5}

Interface conditions for nonlocal media differ from the classical ones because the stress at a point depends on the strain in a neighbourhood, and a neighbourhood straddling an interface samples both media. The classical conditions---continuity of displacement and of traction---must be supplemented, and the supplementation is not unique: one may impose conditions on the additional derivatives that appear in the differential form of Eq.~\eqref{eq:7.2}, or one may impose the conditions in the integral form by requiring the nonlocal stress to be continuous across the interface. The two are not equivalent, and the choice made determines the number and kind of boundary conditions that the resulting boundary-value problem admits. For the differential form, the interface conditions include both the classical conditions and the higher-order conditions arising from the $\nabla^2$ term.

Interface problems in which the nonlocal operator is combined with an additional coupled field have been analysed in the reviewed corpus for a fluid in contact with a nonlocal piezothermoelastic medium \cite{Gupta2023IJNMF} and for nonlocal Kelvin--Voigt piezo-semiconductors under a memory-dependent heat flux \cite{Gupta2026HeatMap}. In both cases the analysis proceeds through the modal machinery of Section~\ref{sec:interfaces}, and the specific difficulty in both is the same: the nonlocal operator acts on the stress while the coupled field acts on the temperature or the electric displacement, and the interface conditions must be formulated consistently for both.

\subsection{What nonlocality does not repair}
\label{sec:7.6}

A review that presented nonlocal elasticity purely as a remedy for the breakdown of classical elasticity would give a misleading impression, because the remedy has a well-documented cost. The differential form of Eq.~\eqref{eq:7.2} raises the order of the governing equation, and the resulting initial-value problem is not well posed in the sense of Hadamard: the number of required initial conditions grows with the order of the equation, and the growth of the solution with the wave number is unbounded \cite{Eringen1983}. In the frequency domain, the same statement appears as the unbounded increase of the wave number with frequency, so that the group velocity grows without bound and the causal cone is not preserved.

This is not a defect of any particular implementation; it is a structural property of the class. The two classical remedies are the exclusion of the problematic high-wave-number branch by a bounded kernel, and the reformulation of the theory in the variational setting, in which the constitutive law appears as a quadratic functional and the resulting evolution problem is well posed in a suitable energy space \cite{Polizzotto2001}. The second remedy is the more satisfactory one, and it is the reason the variational structure of nonlocal elasticity has been worked out with particular care. It also explains the emphasis in Section~\ref{sec:methods} on variational formulations: for nonlocal and gradient theories, the variational route is not an alternative to the differential route but the more rigorous of the two.

\subsection{Comparative summary: nonlocal models}

\begin{table}[htbp]
\centering
\small
\caption{Comparative summary of the nonlocal and gradient models, with the physical mechanism each encodes and the signature by which they could be distinguished experimentally.}
\label{tab:nonlocal_comparison}
\begin{tabularx}{\textwidth}{@{}lXXXX@{}}
\toprule
\textbf{Model} & \textbf{Length scale} & \textbf{Physical mechanism} & \textbf{Dispersion signature} & \textbf{Energy flux form} \\
\midrule
Eringen integral & $e_0 a$ & Neighbourhood averaging of the local stress & Saturating, $c \to 0$ as $k \to \infty$ & Integral form with correction of Eq.~\eqref{eq:7.6} \\
Eringen differential & $\xi = e_0 a$ & Neighbourhood averaging, Helmholtz kernel & Saturating branch, unbounded $k$ & Modified classical form; well-posedness needs care \\
Strain-gradient & $\ell$ & Curvature of the strain field & Linear decrease of $c^2$ in $k^2$ within the valid range & Higher-order gradient terms \\
Couple-stress / micropolar & $\ell$ & Independent micro-rotation field & Depends on the closure of the micropolar equations & Additional micro-rotational work terms \\
Two-phase local/nonlocal & $e_0 a$ plus local weight & Partial averaging of the local stress & Saturating, with a local branch retained & Two-channel flux \\
\bottomrule
\end{tabularx}
\end{table}

\subsection{Notation summary for Section 7}

\begin{table}[htbp]
\centering
\small
\caption{Notation introduced in Section 7.}
\label{tab:notation-sec7}
\begin{tabularx}{\textwidth}{@{}lXlr@{}}
\toprule
\textbf{Symbol} & \textbf{Meaning} & \textbf{Units} & \textbf{Equation} \\
\midrule
$\alpha(|\bm{x}'-\bm{x}|,\tau)$ & Nonlocal attenuation kernel & varies & \eqref{eq:7.1} \\
$\tau$ & Nonlocal parameter, ratio of internal to external length & dimensionless & \eqref{eq:7.1} \\
$e_0$ & Material constant of the nonlocal formulation & dimensionless & \eqref{eq:7.1} \\
$a$ & Internal characteristic length & m & \eqref{eq:7.1} \\
$\xi = e_0 a$ & Nonlocal length scale parameter & m & \eqref{eq:7.2} \\
$\ell$ & Strain-gradient or couple-stress length scale & m & \eqref{eq:7.3} \\
$c(k)$ & Wave speed as a function of wave number & m\,s$^{-1}$ & \eqref{eq:7.4} \\
$c_0$ & Classical (local) wave speed & m\,s$^{-1}$ & \eqref{eq:7.4} \\
$J_i^{\mathrm{nonlocal}}$ & Nonlocal energy flux correction & W\,m$^{-2}$ & \eqref{eq:7.5} \\
\bottomrule
\end{tabularx}
\end{table}

\subsection{Open questions for Section 7}

\begin{enumerate}[label=\textbf{OQ7.\arabic*}]
\item What is the correct energy flux at a nonlocal interface, derived from a rigorous modified divergence theorem, and does it reproduce the classical limit as $\xi \to 0$? \label{oq:7.1}
\item How must the boundary conditions be formulated for the differential nonlocal operator so that the initial-value problem is well posed, and is the resulting formulation variational? \label{oq:7.2}
\item What is the relationship between the integral and differential nonlocal formulations at the level of the energy flux, given that the two agree only for the specific kernel used to obtain the differential form? \label{oq:7.3}
\item Can a dispersion measurement over a sufficiently broad range of wave number distinguish a nonlocal mechanism from a gradient mechanism in a specific material? \label{oq:7.4}
\item Which nonlocal constitutive law is thermodynamically admissible when combined with a memory-dependent thermal law, given that both introduce history into the energy balance? \label{oq:7.5}
\end{enumerate}

% LocalWords:
